\section{서론}

\subsection{동기}

제조공장의 전력 관리는 하루 평균보다 최대수요전력을 먼저 봐야 한다. 공모 과제 ⑤는 공장 전력이 설비 동시 가동, 생산량 변화, 제품 전환, 교대 시간, 계절 요인에 따라 달라지고, 최대수요전력이 전력비용과 생산 안정성에 영향을 준다고 설명한다. 평균 전력은 생산량을 따라 완만하게 움직이지만 최대수요전력은 여러 설비가 같은 15분 안에 겹칠 때 한 번에 올라간다. 부하가 한 시간대에 몰리면 수전 설비의 여유가 줄고 설비 운전도 불안정해진다.

산업용 전기요금 구조에서는 단 하루가 한 달, 나아가 이후 여러 달의 기본요금을 정한다. 기본요금은 요금 시간대(일요일·공휴일을 제외한 09--23시의 중간·최대부하 시간대)의 월 최대수요전력으로 매겨지고, 12·1·2·7·8·9월의 최대수요전력은 이후 달의 하한(래칫)으로도 쓰인다. 본 보고서가 분석한 공장에서는 2021년 7월 19일이 이런 날이었다. 이날 요금 시간대 최대는 222 kW로, 그 전까지의 기준선 206 kW를 16 kW 넘었다. 기본요금 8,320원/kW·월을 적용하면 그 달에만 133,120원이 더 붙고, 7월은 래칫 적용월이므로 이 값은 이후 달의 하한에도 남을 수 있다. 검증 구간(7월 7일--8월 31일)에서 요금 시간대가 있는 44일 가운데 월 최대를 새로 쓴 날은 이 하루뿐이었다.

그래서 예측은 운영 결정이 가능한 시점에, 결정에 쓰는 정보로 내야 한다. 본 보고서의 예측은 매일 00:00에 발행한다. 입력은 당일의 시간별 생산계획, 요일과 공휴일 같은 달력 정보, 전날까지의 전력 관측이며, 출력은 당일 96개 15분 슬롯의 전력 분포다. 생산계획은 공장이 전날에 이미 알고 있고 바꿀 수도 있는 입력이다. 이 정보로 전력을 예측하면 결과를 계획 점검과 연결할 수 있다.

실무에서 필요한 답은 평균 정확도만이 아니다. 운영자는 오늘이 최대 피크 위험이 높은 날인지, 그렇다면 무엇을 점검해야 하는지를 알고 싶어 한다. 슬롯별 오차가 작아도 고피크일의 최댓값을 크게 낮춰 잡는 모형은 요금 관점에서 쓸모가 적다. 반대로 매일 경보를 내면 현장은 곧 경보를 무시한다. 본 보고서는 15분 곡선의 정확도, 일 피크의 분포, 월 최대 갱신 위험을 한 모형에서 함께 산출하고 각각을 따로 채점한다.

\subsection{문헌연구}

\paragraph{공장 부하 예측과 생산 정보.} 생산 정보를 입력하면 공장 전력 예측이 좋아진다는 근거는 쌓였지만, 시간별 생산계획을 15분 전력의 확률분포로 옮긴 연구는 드물다. Walther와 Weigold의 체계적 문헌검토는 제조업 전력 연구 72편 가운데 공장 단위 예측을 다룬 논문이 3편뿐임을 보였다 \citep{walther2021review}. 제지 공장 연구에서는 생산관리 정보를 넣자 MAPE가 평균 18.3\% 개선되었다 \citep{lai2022papermaking}. \citet{trenz2023production}은 생산계획과 기상 예보로 15분 전력을 예측하는 도구를 제시했다. 국내 대형 제조공장을 다룬 \citet{kim2023c2transformer}은 일 작업량으로 일 피크를 먼저 회귀한 뒤 이를 조건으로 시간별 피크 계열을 생성했는데, 시간별 계획은 실무에서 드물다고 보고 일 단위 작업량만 썼다. 조선소 사례 \citep{lee2026shipyard}는 공정 지표·달력·기상을 입력으로 일 피크를 예측하는 12개 모형을 비교하고, 래칫 요금 아래서는 하루의 높은 피크가 연간 비용을 좌우한다고 지적했다. \citet{berk2018regime}은 산업 부하의 가동과 정지를 레짐 전환으로 보고 확률 예측을 했다. 이 연구들은 대부분 점예측을 MAPE·MAE로 평가하거나 일 단위 피크만 다룬다.

\paragraph{확률 예측과 피크 위험.} 확률 예측의 평가 체계는 자리를 잡았지만 하루 최댓값은 시점별 분위수에서 바로 얻을 수 없다. \citet{hong2016tutorial}은 확률적 부하 예측의 입력, 모형, 평가를 정리했고, GEFCom2014는 분위수와 pinball loss를 표준 평가로 정착시켰다 \citep{hong2016gefcom}. pinball loss는 분위수 회귀 \citep{koenker1978quantiles}의 점검 함수이며, CRPS와 함께 엄밀히 적절한 점수 규칙에 속한다 \citep{gneiting2007scoring}. 피크를 확률로 다룬 연구는 두 갈래로 나뉜다. 하나는 피크를 직접 예측하는 방식이다. \citet{hyndman2010peak}은 연간 피크의 밀도와 초과확률을 제시했고, \citet{amaraouali2023peak}은 고해상도와 저해상도 정보를 함께 써서 일 피크의 크기와 시각을 예측했다. \citet{gilbert2023peaks}은 일반 예측이 지나치게 매끄러워 피크를 낮춰 잡는다는 점을 지적하고 일 피크 예측을 30분 예측과 융합해 피크 시간대 CRPS를 개선했다. 다른 하나는 하루 전체 경로의 공동분포를 표집한 뒤 그 최댓값을 취하는 방식이다. 이 방식은 곡선, 피크, 초과확률을 하나의 분포에서 일관되게 얻지만 경로 잡음의 크기가 맞지 않으면 피크가 한쪽으로 치우친다. 본 보고서는 경로 방식을 택하고, 이 치우침을 가동 상태별 잡음으로 다룬다.

\paragraph{베이지안 평활 모형과 유사일.} 곡선의 형태, 입력의 시차 효과, 상태에 따른 분산은 각각 확립된 도구가 있지만 한 사후분포로 묶여 쓰인 예는 찾지 못했다. RW2 사전분포는 가우스 마르코프 확률장의 대표 사례로 곡선을 매끄럽게 추정한다 \citep{rue2005gmrf}. 베이지안 분포시차모형은 시차 계수에 평활 제약을 두어 입력이 여러 시차에 걸쳐 전달되는 구조를 추정한다 \citep{welty2009bdlm}. GAMLSS는 분산 같은 분포 모수를 공변량의 함수로 둔다 \citep{rigby2005gamlss}. 과거 사례를 빌리는 쪽에서는 유사일(analog) 예측이 오래 쓰였고, \citet{mcdermott2016analog}은 이를 베이지안 모형으로 정식화해 예측 불확실성을 정량화했다. 학습형 어텐션도 과거 관측과 미래에 알려진 입력을 결합하는 방법이다 \citep{lim2021tft}. sparsemax는 가중치를 정확히 0으로 만들 수 있어 소수의 참조만 남기는 희소한 어텐션을 준다 \citep{martins2016sparsemax}. 그러나 학습형 가중치는 작은 공장 자료에서 불안정하다. 본 보고서는 sparsemax를 고정된 국소 커널로 써서 분포시차 전달함수처럼 해석한다.

\paragraph{수요요금과 생산 일정.} 피크 저감 스케줄링 연구는 설비 전력을 알려진 값으로 두므로 예측 오차가 결정에 미치는 영향을 다루지 않는다. \citet{gahm2016ees}은 에너지 효율 스케줄링 연구를 정리했고, \citet{fang2011peakpower}은 피크 전력을 목적함수에 넣은 flow shop 모형을 제시했다. \citet{shrouf2014tou}은 시간대별 요금(TOU)에서 설비 가동 시점을 최적화했고, \citet{shoreh2016dr}은 산업 수요반응의 잠재력과 도입 장벽을 조사했다. 이 흐름은 계획을 바꾸면 피크가 얼마나 줄지를 결정론적으로 계산한다. 예측 연구는 반대로 조치 단계까지 가지 않는다.

네 갈래를 합치면 공백이 드러난다. 시간별 생산계획을 조건으로 15분 전력의 공동 경로를 만들고, 그 경로에서 일 피크 분포와 월 최대 갱신 위험을 함께 산출해 요금 규칙에 맞춰 채점한 공장 단위 연구는 찾지 못했다. 본 보고서는 이 공백을 다룬다.

\subsection{기여}

첫째, 생산계획을 조건으로 한 베이지안 전달모형 C-BAT를 제안하고 개발 구간과 9월을 합친 시간순 평가에서 피크 예측을 개선했다. C-BAT는 가동유형별 기준 곡선, 최근 기준일, 생산계획의 양(+)의 전달항을 결합하고 잡음의 크기를 가동 상태에 따라 다르게 둔다. 합산 평가는 67일·6,358슬롯이며 피크 평가일은 65일이다. 피크 MAE는 C-BAT 8.93 kW, 기존 BAT 13.09 kW, LightGBM 기본 모형 M2 14.54 kW로, C-BAT가 각각 31.8\%, 38.6\% 낮았다. 날짜 단위 짝지은 bootstrap의 차이는 기존 BAT 대비 $-4.17$ kW(95\% 구간 $[-6.95,-1.31]$), M2 대비 $-5.61$ kW($[-10.24,-1.38]$)였다. 슬롯 MAE는 C-BAT 8.22 kW, 기존 BAT 7.61 kW, M2 10.57 kW로 피크와 순위가 다르다. 이 합산 비교에는 B0\_kind와 튜닝 LightGBM이 포함되지 않았다.

개발 구간 f1--f4에서는 더 많은 기준선과 조건별 성능을 비교했다(53일, 피크 평가 51일). 고피크일 13일의 피크 MAE는 C-BAT 4.37 kW로, 튜닝 LightGBM \citep{ke2017lightgbm}의 18.27 kW와 기존 BAT의 12.47 kW보다 작았다. 튜닝 LightGBM과의 차이는 $-13.90$ kW(날짜 bootstrap 95\% 구간 $[-19.01,-9.86]$)였다. 반면 개발 구간 전체 피크 MAE는 기준일 복사 기준선 B0\_kind가 6.67 kW로 C-BAT(8.93 kW)보다 낮았다. 모형 선택은 이 개발 구간에서 끝내고 9월은 합산 성능 보고에만 사용했다.

둘째, 사후 경로로 월 최대 갱신 위험을 계산하는 경보를 만들고 과거 날짜 백테스트로 채점했다. 경보는 요금 시간대 경로 최댓값이 그날까지의 기준선을 넘는 정도로 기대 초과 기본요금을 계산하고, 이 값이 3만 원 이상이면 점검을 권한다. 검증 구간의 유일한 갱신일인 7월 19일에 C-BAT는 갱신 확률 0.58, 기대 초과 비용 122,919원을 내어 경보했고, B0\_kind는 195 kW를 예측해 놓쳤다. 주 규칙의 경보는 44일 중 10회였고 그중 9회는 오경보였다. 갱신일이 하루뿐이므로 이 결과는 성능 추정이 아니라 사례로 제시한다.

셋째, 자료의 복사일과 평가 누수를 먼저 진단하고 채택 규칙을 결과 확인 전에 고정했다. 첫 검증 구간의 학습기간 186일 가운데 121일이 전력 곡선이 복사된 날이었다. 복사일이 섞인 자료에서 날짜를 무작위로 나누면 LightGBM의 MAE가 14.91 kW로, 복사 사건 단위로 떼어 낸 평가의 16.64 kW보다 약 10\% 좋게 보였다 \citep{roberts2017blockcv,kapoor2023leakage}. 그래서 모든 비교는 복사일을 뺀 시간순 분할로 했다. 확장 실험(계획 기반 기준일, 어텐션 기준일, 피크 계층, 앙상블)은 사전 등록한 규칙으로 판정했고 채택된 것은 없다. 9월 1--14일 구간에서는 C-BAT의 슬롯 MAE가 7.18 kW, 피크 MAE가 8.93 kW로 LightGBM 기반 M2(6.60 kW, 5.53 kW)보다 나빴다. 이 구간은 이전 개발 중에 한 번 열린 적이 있어 모형 선택 근거로 쓰지 않았지만, 결과는 그대로 보고한다.

넷째, 현재 생산계획과 시도한 모형만으로 평범한 가동일과 고피크일을 충분히 구별하기 어렵다는 한계를 확인했다. 계획 기반 기준일은 슬롯 MAE를 0.87 kW 낮췄지만 가동일 피크 과대예측을 바꾸지 못했다. 기준일 피크를 가동유형별로 재보정한 피크 계층은 두 설계 모두 고피크일을 20 kW 이상 낮춰 잡았다. C-BAT 경보가 놓친 고피크일도 모두 7월의 가장 긴 가동유형에 몰렸다. 이 결과는 생산계획에 정보가 전혀 없다는 증명이 아니라, 설비별 가동 일정이나 제품 전환 기록을 추가해 구별력을 다시 검증할 근거다.

2장은 C-BAT의 구조와 경로 기반 피크 위험 계산을 설명한다. 3장은 자료의 복사일, 결측, 피크 발생 조건을 분석한다. 4장은 비교 모형과의 결과와 오류 조건을 제시한다. 5장은 월 최대 갱신 경보의 현장 활용방안을, 6장은 차별성과 결론을 다룬다. 부록 A는 재현 절차를 담는다.
